\documentclass{amsart}
\numberwithin{equation}{section}

%%%%%%%packages
\usepackage{amssymb}
\usepackage{enumerate}
%\usepackage{showkeys}
\usepackage[bottom,first]{draftcopy}
\usepackage{xspace}
\usepackage{hyperref}

%%%%%%%%%Pagination
\hfuzz=15pt

%%%%%%%%%Theorems
\newtheorem{thm}{Theorem}[section]
\newtheorem{lem}[thm]{Lemma}
\newtheorem{cor}[thm]{Corollary} 
\newtheorem{sublem}[thm]{Sub-lemma}
\newtheorem{prop}[thm]{Proposition}
\newtheorem{conj}{Congettura}\renewcommand{\theconj}{}
\newtheorem{defin}[thm]{Definition}
\newtheorem{cond}{Condition}
\newtheorem{exmp}[thm]{Example}
\newtheorem{es}[thm]{Exercize}
\newtheorem{rem}[thm]{Comment}

%%%%%%%%%%mathcal
\newcommand\B{{\mathcal B}}
\newcommand\Co{{\mathcal C}}
\newcommand\D{{\mathcal D}}
\newcommand\F{{\mathcal F}}
\newcommand\Lp{{\mathcal L}}
\newcommand\M{{\mathcal M}}
\newcommand\Or{{\mathcal O}}
\newcommand\Pa{{\mathcal P}}
\newcommand\T{{\mathcal T}}

%%%%%%%%%%%%mathbb
\newcommand\A{{\mathbb A}}
\newcommand\C{{\mathbb C}}
\newcommand\E{{\mathbb E}}
\newcommand\N{{\mathbb N}}
\newcommand\R{{\mathbb R}}
\newcommand\Q{{\mathbb Q}}
\newcommand\To{{\mathbb T}}
\newcommand\V{{\mathbb V}}
\newcommand\Z{{\mathbb Z}}

%%%%%%%%%%%greek
\newcommand\ve{\varepsilon}
\newcommand\vf{\varphi}

%%%%%%%%%%%%%%%vaious
\newcommand\Id{\text{\bf Id}}

%%%%%   Figure-picture related definitions--beginning
\newenvironment{myfigure}{\begin{figure}[tb]\
\centering}{\end{figure}}

%\definecolor{lgray}{gray}{0.9}
%\definecolor{gray}{gray}{0.8}
%\definecolor{dgray}{gray}{0.7}
\setlength{\unitlength}{.8mm}
%%%%%   Figure related definitions--end


\pagestyle{myheadings}
\markboth{\centerline{\sc Carlangelo Liverani}\hskip-12pt}
{\centerline{\sc Why Dymanics matters? }\hskip-12pt}

\begin{document}
\title{\bfseries A piece of functional analysis}
\author{Carlangelo Liverani}
\address{Dipartimento di Matematica\\
Universit\`{a} di Roma (Tor Vergata)\\
Via della Ricerca Scientifica, 00133 Roma, Italy\\
{\tt liverani@mat.uniroma2.it}}
\date{Rome, October 11, 2003}
\maketitle


\section{The setting}\label{sec:zero}

Consider two Banach space $(\B,\|\cdot\|)$ and $(\B_0,|\cdot|)$ such
that $\B\subset \B_0$ and
\begin{enumerate}[i.\hskip.4cm]
\item $|h|\leq \|h\|$ for all $h\in\B$,
\item if $h\in\B$ and $|h|=0$, then $h=0$.
\item The unit ball $\{h\in\B\;:\;\|h\|\leq 1\}$ is relatively compact in $\B_0$.
\item There exists $C>0$ such that: for each $\ve>0$ there
exists a finite rank operator $\A_\ve\in L(\B,\B)$ such that $\|\A_\ve\|\leq
C$ and $|h-\A_\ve h|\leq \ve \|h\|$ for all $h\in\B$.\footnote{In fact,
this last property is not needed. We require it since, on the one
hand, it is true in all the 
applications we will be interested in and, on the other hand,
drastically simplifies the argument. Note also that, if one uses the
Fredholm alternative for compact operators rather than finite rank
ones (Theorem \ref{thm:fredholm}), then one can ask the $\A_\ve$ to be compact
instead than finite rank making easier their construction in concrete cases.}
\end{enumerate}

In addition consider a bounded operator $\Lp:\B_0\to\B_0$,
constants $A,B,C\in\R_+$, and $\lambda>1$, such that
\begin{enumerate}[a.\hskip.4cm]
\item   $|\Lp^n|\leq C$ for all $n\in\N$,
\item  $\Lp(B)\subset B$ 
\item  $\|\Lp^nh\|\leq A\lambda^{-n}\|h\|+B|h|$ for all
$h\in\B$ and $n\in\N$. 
\end{enumerate}
In particular $\Lp$ can be seen as a bounded operator on $\B$.

\begin{thm}\label{thm:main}
The spectral radius of the operator $\Lp\in L(\B,\B)$ is bounded by
$1$ while the essential spectral radius is bounded by $\lambda^{-1}$.
\end{thm}

The proof of the above theorem depends on an auxiliary fact which is of
much interest in its own.

\begin{thm}[Analytic Fredholm theorem--finite rank\footnote{The present proof is
patterned after the proof of the Analytic Fredholm alternative for
compact operators (in Hilbert spaces) given in \cite[Theorem
VI.14]{RS}. There it is used the fact that compact operators in
Hilbert spaces can always be approximated by finite rank ones. In fact
the theorem holds also for compact operators in Banach spaces but the
proof is a bit more involved.}]\label{thm:fredholm}
Let $D$ be an open connected subset of $\C$. Let $F:\C\to L(\B,\B)$ be
an analytic operator-valued function such that $F(z)$ is finite rank
for each $z\in D$. Then, one of the following two alternatives holds true
\begin{itemize}
\item $(\Id-F(z))^{-1}$ exists for no $z\in D$
\item $(\Id-F(z))^{-1}$ exists for all $z\in D\backslash S$ where $S$
is a discrete subset of $D$ (i.e. $S$ has no limit points in $D$). In addition, if
$z\in S$, then $1$ is an eigenvalue for 
$F(z)$ and the associated eigenspace has finite multiplicity.
\end{itemize}
\end{thm}
\begin{proof}
First of all notice that, for each  $z_0\in D$ there exists
$r>0$ such that $D_{r(z_0)}(z_0):=\{z\in\C\;:\;
|z-z_0|<r(z_0)\}\subset D$, and
\[
\sup_{z\in D_{r(z_0)}(z_0)}\|F(z)-F(z_0)\|\leq \frac 12.
\]
Clearly if we can prove the theorem in each such disk we are
done.\footnote{In fact, consider any connected compact set $K$
contained in $D$. Let us suppose that
for each $z_0\in K$ we have a disk $D_{r(z_0)}(z_0)$ in the theorem
holds.
Since the disks $D_{r(z_0)/2}(z_0)$ form a covering for $K$ we can
extract a finite cover. If the first alternative holds in one such
disk then, by connectness, it must hold on all $K$. Otherwise each
$S\cap D_{r(z_0)/2}(z_0)$, and hence $K\cap S$, contains only finitely
many points. The Theorem follows by the
arbitrariness of $K$.}
Note that
\[
\Id-F(z)=\left(\Id-F(z_0)(\Id-[F(z)-F(z_0)])^{-1}\right)(\Id-[F(z)-F(z_0)]).
\]
Thus the invertibility of $\Id-F(z)$ in $D_r(z_0)$ depends on the
invertibility of $\Id-F(z_0)(\Id-[F(z)-F(z_0)])^{-1}$. Let us set
$F_0(z):=F(z_0)(\Id-[F(z)-F(z_0)])^{-1}$.

Let us start by looking at the equation
\begin{equation}\label{eq:ker}
(\Id-F_0(z))h=0.
\end{equation}
Clearly if a solution exists, then
$h\in\text{Range}(F_0(z))=\text{Range}(F(z_0)):=\V_0$. Since $\V_0$ is
finite dimensional  
there exists a basis $\{h_i\}_{i=1}^N$ such that
$h=\sum_i\alpha_i h_i$. On the other hand there exists an analytic
matrix $G(z)$ such that\footnote{To see the analyticity notice that we
can construct linear functionals $\{\ell_i\}$ on $\V_0$ such that
$\ell_i(h_j)=\delta_{ij}$ and then extend them to all $\B$ by the
Hahn-Banach theorem. Accordingly, $G(z)_{ij}:=\ell_j(F_0(z)h_i)$, which
is obviously analytic.}
\[
F_0(z)h=\sum_{ij}G(z)_{ij}\alpha_j h_i.
\]
Thus \eqref{eq:ker} is equivalent to
\[
(\Id-G(z))\alpha=0,
\]
where $\alpha:=(\alpha_i)$.

The above equation can be satisfied only if $\det(\Id-G(z))=0$ but the
determinant is analytic hence it is either always zero or zero only at
isolated points.\footnote{The attentive reader has certainly noticed that this
is the turning point of the theorem: the discreteness of $S$ is
reduced to the discreteness of the zeroes of an appropriate analytic
function: a determinant. A moment thought will immediately explain the
effort made by many mathematicians to extend the notion of determinant
(that is to define an analytic function whose zeroes coincide with the
spectrum of the operator) 
beyond the realm of matrices (the so called Fredholm determinants).}

Suppose the determinant different from zero, and consider the equation
\[
(\Id-F_0(z))h=g.
\]
Let us look for a solution of the type $h=\sum_i\alpha_i h_i
+g$. Substituting yields
\[
\alpha- G(z)\alpha=\beta
\]
where $\beta:=(\beta_i)$ with $F_0(z)g=:\sum_i\beta_i h_i$. Since the
above equation admits a 
solution, we have $\text{Range}(\Id-F_0(z))=\B$, Thus we have an
everywhere defined inverse, hence bounded by the open mapping theorem.

We are thus left with the analysis of the situation $z\in S$ in the
second alternative. In such a case, there exists $h$ such that
$(\Id-F(z))h=0$, thus one is an eigenvalue. On the other hand, if we
apply the above facts to the function $\Phi(\zeta):=\zeta^{-1}F(z)$
analytic in the domain $\{\zeta\neq 0\}$ we note that the first
alternative cannot take place since for $|\zeta|$ large enough
$\Id-\Phi(\zeta)$ is obviously invertible. Hence, the spectrum of
$F(z)$ is discrete and can accumulate only at zero. This means that
there is a small neighborhood around one in which $F(z)$ has no
other eigenvalues, we can thus surround one with a small circle
$\gamma$ and consider the projector
\[
\begin{split}
P:=&\frac1{2\pi i}\int_\gamma (\zeta-F(z))^{-1}d\zeta=\frac1{2\pi
i}\int_\gamma \left[(\zeta-F(z))^{-1}-\zeta^{-1}\right]d \zeta\\
=&\frac1{2\pi i}\int_\gamma F(z)\zeta^{-1}(\zeta-F(z))^{-1}d\zeta.
\end{split}
\]
By standard functional calculus it follows that $P$ is a projector and
it clearly projects on the eigenspace of the eigenvector one. But
the last formula shows that $P$ must project on a subspace of the rank
of $F(z)$, hence it must be finite dimensional.
\end{proof}

We can now prove our main result.

\begin{proof}[Proof of Theorem \ref{thm:main}]
The first assertion is a trivial consequence of (c), (a) and (i). 

The second part is much deeper.
Let $\Lp_{n,\ve}:=\Lp^n\A_\ve$, clearly
such an operator is finite rank, in addition
\[
\|\Lp^n h-\Lp_{n,\ve}h\|\leq
A\lambda^{-n}\|(\Id-\A_\ve)h\|+B|(\Id-\A_\ve)h|
\leq A(1+C)\lambda^{-n}\|h\|+B\ve\|h\|.
\]
By choosing $\ve=\lambda^{-n}$ we have that there exists $C_1>0$ such
that
\[
\|\Lp^n -\Lp_{n,\ve}\|\leq C_1\lambda^{-n}.
\]
For each $z_0\in\C$ we can now write
\[
\Id-z\Lp=(\Id-z_0(\Lp-\Lp_{n,\ve})-(z-z_0)\Lp)-z_0\Lp_{n,\ve}.
\]
Since
\[
\|z_0(\Lp-\Lp_{n,\ve})+(z-z_0)\Lp\|\leq
|z_0|C_1\lambda^{-n}+C|z-z_0|<\frac 12,
\]
provided that $|z_0|\leq \frac 1{4C_1}\lambda^{n}$ and $|z-z_0|\leq
\frac 1{4C}$, given any $z_0$ in the disk $D_n:=\{|z|<\frac
1{4C_1}\lambda^{n}-\frac 1{4C}\}$ we can consider the domain
$D_n(z_0):=\{|z-z_0|<\frac 1{4C}\}$ and, in such a domain
$B(z):=\Id-z_0(\Lp-\Lp_{n,\ve})-(z-z_0)\Lp$ is invertible. Hence
\[
\Id-z\Lp=
\left(\Id-
z_0\Lp_{n,\ve}B(z)^{-1}\right)B(z)=:
(\Id-F(z))B(z).
\]
By applying Theorem \ref{thm:fredholm} to $F(z)$ we have that the
operator is either never invertible or not invertible only in finitely
many points in the disk $D_n'(z_0):=\{|z-z_0|<\frac 1{8C}\}$. Since we
can cover $D_n:=\{|z|<\frac
1{4C_1}\lambda^{n}-\frac 1{4C}\}$ with finitely many such disk and
since the operator is invertible for $z_0$ small enough ( thus the
first alternative cannot hold), the Theorem follows.
\end{proof}


\begin{thebibliography}{99}
\footnotesize
\bibitem{RS} M.Reed, B.Simon, {\em Methods of Modern Mathematical
Physics. I-Functional Analysis}, Academic Press, New York (1980). 
\end{thebibliography}
\end{document}
